\documentclass[acmsmall,screen,review,anonymous,nonacm]{acmart}

\usepackage{amsmath,amsthm}
\usepackage{array}
\usepackage{mathtools}
\usepackage{listings}

\settopmatter{printfolios=true,printccs=false,printacmref=false}
\citestyle{acmnumeric}

\lstset{
  basicstyle=\ttfamily\small,
  columns=fullflexible,
  keepspaces=true,
  showstringspaces=false,
  frame=single
}

\newtheorem{definition}{Definition}
\newtheorem{proposition}{Proposition}
\newtheorem{theorem}{Theorem}

\newcommand{\Types}{\mathcal{T}}
\newcommand{\Effects}{\mathcal{E}}
\newcommand{\glb}{\mathbin{\sqcap}}
\newcommand{\compose}{\mathbin{\cdot}}
\newcommand{\eps}{\varepsilon}
\newcommand{\id}{\mathbf{1}}
\newcommand{\sem}[1]{[\![#1]\!]}

\title{Declarative Stack-Effect Checking for a Forth-like Language}
\author{Anonymous Author(s)}
\renewcommand{\shortauthors}{Anonymous Author(s)}

\begin{document}

\begin{abstract}
This paper presents a static checker for a Forth(-like) language based on typed stack effects. The checker is parameterized by a finite type hierarchy together with declarative specifications for ordinary words, literals, defining words, parser words, and control constructs. Its core operations are sequential composition of effects, branch merging by greatest lower bound in a must-analysis order, and loop summarization by a local stabilization criterion. Control forms are described by paired \texttt{syntax:}/\texttt{effect:} declarations, so source-level parsing and effect computation are expressed in one framework. For a fixed parsed source tree whose leaves already have declared summaries, checking admits an exact compositional characterization: each parsed segment receives a unique abstract summary obtained by recursively evaluating leaf effects, sequential composition, and declarative control templates, and acceptance is exactly the absence of clashes in that evaluation. The implementation, written in Gforth, supports conditionals, loops, and user-defined words. We present the effect calculus, the compositional theorem, and an evaluation on a 29-program regression corpus, retargetable custom profiles, and a 4{,}255-line application.
\end{abstract}

\maketitle

\section{Introduction}
Concatenative languages compose programs by placing words next to one another, so juxtaposition itself denotes composition. In stack-based instances such as Forth, words perform computation by transforming an implicit operand stack. A word is therefore naturally described by the stack cells it consumes and the cells it produces, so a natural static abstraction is a \emph{stack effect}. In simple cases, informal comments such as \texttt{( n n -- n )} are already useful documentation in Forth practice and standardization documents \cite{rather1996evolution,forth2012standard}. The goal of the checker described here is to turn such annotations into an executable static semantics.

The checker takes three external inputs. A \texttt{types} file defines a finite subtype hierarchy. A \texttt{specs} file assigns effects and parsing behavior to words, literals, defining words, and control roles. A program file is then parsed and analyzed against those specifications. Instead of running the program for its concrete result, the checker computes the abstract stack effect of each token and of each user-defined word. The same mechanism produces source-aware diagnostics when adjacent effects cannot be composed. In that sense, the implementation follows the abstract-interpretation idea of executing programs over a finite abstract domain rather than over concrete values \cite{cousot1977abstract}.

A central technical result is that this operational account has an exact declarative reading. Once parsing has segmented a source body and the referenced words have fixed summaries, every parsed segment is assigned a unique abstract effect by a finite family of equations over the source tree. The theorem proved later states that the checker accepts exactly when these equations evaluate without clash, and that on success the reported summaries are precisely the induced ones. This gives a compositional characterization of the checking procedure that matches the implementation more closely than a global fixpoint story.

Two design choices make the system concise and extensible. First, types are not hard-coded; they are loaded as a small finite lattice-like structure. Second, control flow is not built into the checker alone; it can be described declaratively with \texttt{syntax:} and \texttt{effect:} blocks. The mathematical content of the checker is therefore concentrated in a small abstract domain together with a few monotone effect transformers. This approach is informed by earlier work on algebraic specifications of stack effects, typed wildcards, must-analysis for stack-oriented languages, and compositional type systems for stack-oriented languages \cite{poial1994algebraic,stoddart1993type,poial2002wildcards,saabas2006compositional,poial2006typing}.

Earlier papers in this line developed the ingredients separately: algebraic composition of stack effects \cite{poial1994algebraic}, typed wildcards and subtype-sensitive symbolic refinement \cite{stoddart1993type,poial2002wildcards}, and must-analysis formulations based on greatest lower bounds together with loop summaries extracted from repeated effects \cite{poial2006typing}. The present paper is concerned with combining these ingredients into one account for whole source programs and relating that account to the older inequality-based view of stack-effect specifications \cite{poial1991algebraic,poial1992multiple}.

Concretely, the paper makes four contributions.
\begin{enumerate}
\item It presents a declarative effect calculus in which stack symbols carry both subtype information and identity constraints, so stack shuffles can be reasoned about compositionally.
\item It gives a must-analysis interpretation of branching and looping in the same effect domain: branches merge by greatest lower bound, while loops are summarized by a lightweight local stabilization test.
\item It proves an exact compositional theorem for fixed parsed trees: the checker accepts exactly when the induced segment equations are clash-free, and on success returns the unique family of summaries satisfying them.
\item It evaluates the calculus as a retargetable checker: the type hierarchy, vocabulary, parser behavior, and control structures are all loaded from external specifications, user-defined words are inferred during analysis, and the artifact is exercised on a 29-program regression corpus, several custom profiles, and a 4{,}255-line application.
\end{enumerate}

Sections~2--5 define the abstract domain and its transfer operations. Section~6 describes declarative control specifications and the compositional characterization. The remaining sections discuss inference of user-defined words, relation to prior work, evaluation, and conclusion.

\section{Type Domain}
\subsection{Type hierarchy}
Let $\Types$ be a finite set of type names equipped with a precision order $\leq$ \cite{davey2002lattices}. We read $\tau_1 \leq \tau_2$ as ``$\tau_1$ is at least as precise as $\tau_2$,'' matching the sample declarations
\begin{lstlisting}
rel a-addr < c-addr
rel c-addr < addr
rel addr < x
rel flag < x
rel char < n
rel n < x
\end{lstlisting}
from the repository. Thus \texttt{a-addr} is more precise than \texttt{c-addr}, which is more precise than \texttt{addr}, and so on.

The implementation stores this relation as a finite transitive closure. Operationally, the checker only needs the meet of \emph{comparable} types.

\begin{definition}[Comparable meet]
For $\tau,\upsilon \in \Types$, define
\[
\tau \glb \upsilon =
\begin{cases}
\tau & \text{if } \tau \leq \upsilon,\\
\upsilon & \text{if } \upsilon \leq \tau,\\
\text{undefined} & \text{otherwise.}
\end{cases}
\]
\end{definition}

When the order is viewed as a finite lattice, this operation coincides with the lattice meet on comparable pairs \cite{davey2002lattices}. Incomparable pairs have no defined result here and therefore cause a type clash.

\subsection{Stack symbols}
A \emph{stack symbol} is a type optionally annotated with a positive index, written $\tau$ or $\tau[i]$. The type denotes an abstract class of stack cells. The index denotes correlation across positions: if the same indexed symbol appears twice, those occurrences refer to the same abstract value constraint, not merely two values of the same coarse type.

For example, the bundled specification
\begin{lstlisting}
PLUS ( x[1] x[1] -- x[1] )
\end{lstlisting}
states that both arguments to \texttt{PLUS} and its result are tied to the same abstract type variable. Likewise,
\begin{lstlisting}
DUP ( x[1] -- x[1] x[1] )
\end{lstlisting}
states that duplication preserves the abstract identity of the top stack element.

Unnamed symbols are treated as implicitly fresh. After evaluation, the checker renumbers indices into a compact normal form, so effects are compared modulo alpha-renaming.

\section{Stack Effects}
\subsection{Syntax}
Let $\Sigma$ be the set of stack symbols. A stack effect is a pair of finite sequences
\[
\pi = (L \; -- \; R), \qquad L,R \in \Sigma^\ast.
\]
The sequence is written in ordinary Forth order: the rightmost element is the top of the stack. For example,
\[
(x[2]\;x[1]\; -- \; x[2]\;x[1]\;x[2])
\]
is the effect of \texttt{OVER}.

The empty effect
\[
\id = (\, -- \,)
\]
acts as the identity transformation.

\subsection{Normalization}
Two effects that differ only by a renaming of fresh indices should be considered equal. The implementation therefore normalizes inferred effects after every successful sequence evaluation. Symbols that occur only once and were never explicitly indexed are printed without an index; correlated or explicitly named symbols are assigned compact indices in a deterministic order.

This normalization step is not cosmetic. It lets the checker treat effects as abstract objects rather than artifacts of a particular evaluation order.

\begin{proposition}[Normalization invariance]
If $\pi'$ is obtained from $\pi$ by alpha-renaming fresh indices, or by the implementation's normalization step, then $\gamma(\pi')=\gamma(\pi)$.
\end{proposition}

\begin{proof}[Proof sketch]
Fresh indices matter only through equality constraints induced by repeated occurrences. Renaming them therefore does not change which concrete stacks admit a common valuation. Likewise, a singleton fresh index contributes no equality constraint at all, so suppressing it in normalized output leaves the concretization unchanged.
\end{proof}

\subsection{Information order}
The branch operator used later is best understood as a meet in a must-analysis order. To make this precise, fix an infinite set $\mathcal{A}$ of concrete atoms. A \emph{concrete stack cell} is a pair $(a,\tau)$ with $a \in \mathcal{A}$ and $\tau \in \Types$; we read $\tau$ as the cell's exact type. A \emph{concrete stack} is a finite sequence of such cells.

Let $\operatorname{Ind}(\pi)$ be the set of indices that occur in an effect $\pi$. A valuation $\eta$ for $\pi$ is a map from $\operatorname{Ind}(\pi)$ to concrete stack cells. If $u$ is a stack symbol with abstract type $\tau_u$, write
\[
c \models_\eta u
\]
when $c=(a,\tau_c)$ satisfies $\tau_c \leq \tau_u$, and if $u$ carries an index $i$, then additionally $\eta(i)=c$. For a symbol sequence $S=u_1\cdots u_n$ and a concrete stack $s=c_1\cdots c_n$ of the same length, write
\[
s \models_\eta S
\]
iff $c_j \models_\eta u_j$ for every position $j$.

For an abstract effect $\pi=(L -- R)$, define its concretization by
\[
\gamma(\pi)=\{(s,t) \mid \text{there exists a valuation } \eta
\text{ such that } s \models_\eta L \text{ and } t \models_\eta R\}.
\]
Thus $\gamma(\pi)$ is the relation on concrete input/output stacks described by $\pi$. Repeated indices induce equality constraints because the same valuation $\eta$ is used on both sides. Unnamed symbols contribute only type constraints, so renaming fresh indices leaves $\gamma$ unchanged.

Define
\[
\pi_1 \preceq \pi_2
\quad\text{iff}\quad
\gamma(\pi_1) \subseteq \gamma(\pi_2).
\]
Lower elements in this order are more informative: they rule out more concrete behaviors and therefore provide stronger guarantees. This is the order-theoretic viewpoint of abstract interpretation specialized to symbolic stack transformations \cite{cousot1977abstract,poial2006typing}.

Under this reading, sequential composition propagates guaranteed information through a linear segment, while branch merging must compute the strongest effect guaranteed by every alternative rather than the union of effects that may arise on some path.

\section{Sequential Composition}
The primary operation of the checker is sequential composition of effects.

\begin{definition}[Boundary composition]
Let
\[
\pi_1=(L_1 -- R_1), \qquad \pi_2=(L_2 -- R_2).
\]
Their composition $\pi_1 \compose \pi_2$ is computed from the touching boundary between the output stack $R_1$ and the input stack $L_2$:
\begin{enumerate}
\item If $R_1$ is empty, the result is $(L_2L_1 -- R_2)$.
\item If $L_2$ is empty, the result is $(L_1 -- R_1R_2)$.
\item Otherwise let $a$ be the top element of $R_1$ and $b$ the top element of $L_2$.
  If $a.\mathsf{type} \glb b.\mathsf{type}$ is undefined, composition fails.
  If it is defined, replace both symbols everywhere by a fresh symbol whose type is that meet, remove the matched boundary pair, and continue recursively.
\end{enumerate}
\end{definition}

The fresh symbol keeps the more specific of the two type constraints. Since lower elements in the order are more precise, the meet is the more informative symbol. This is why an output proved as \texttt{a-addr} can still satisfy a later input that only requires \texttt{c-addr}: the joined boundary keeps the stronger fact \texttt{a-addr}.

Write relational composition of concrete stack transformations as
\[
R_1 ; R_2 = \{(s,u) \mid \exists t.\ (s,t)\in R_1 \text{ and } (t,u)\in R_2\}.
\]

\begin{proposition}[Compositionality of boundary composition]
If $\pi_1 \compose \pi_2$ is defined, then
\[
\gamma(\pi_1 \compose \pi_2) = \gamma(\pi_1) ; \gamma(\pi_2).
\]
\end{proposition}

\begin{proof}[Proof sketch]
The recursive boundary algorithm does exactly the matching needed to eliminate an existentially quantified intermediate stack. At each matched boundary position it replaces the two adjacent symbols by their comparable meet, which is precisely the strongest type constraint satisfied by every concrete intermediate cell that could serve both sides. Shared indices are propagated by the same substitution on both effects, so equality constraints across the boundary are preserved as well. Recursing over the boundary therefore removes the intermediate stack one cell at a time without changing the represented input/output relation.
\end{proof}

The command-line example from the repository illustrates the rule:
\begin{lstlisting}
1 2 + DUP
\end{lstlisting}
with specifications
\begin{lstlisting}
literal integer ( -- n )
+ ( n n -- n )
DUP ( x[1] -- x[1] x[1] )
\end{lstlisting}
The checker prints
\begin{lstlisting}
1   ( --  n[2] )
2   ( --  n[2] )
+   ( n[2] n[2] --  n[1] )
DUP ( n[1] --  n[1] n[1] )
\end{lstlisting}
and concludes that the whole program has effect
\[
( -- n[1]\;n[1]).
\]

\begin{definition}[Multiplication of linear effects]
Let $\Pi_{\mathsf{lin}}$ be the set of normalized linear stack effects and let
\[
\widehat{\Pi}_{\mathsf{lin}} = \Pi_{\mathsf{lin}} \cup \{0\},
\]
where $0$ is a clash element. Define multiplication $\odot$ on $\widehat{\Pi}_{\mathsf{lin}}$ by
\[
0 \odot \pi = \pi \odot 0 = 0,
\]
and for $\pi_1,\pi_2 \in \Pi_{\mathsf{lin}}$,
\[
\pi_1 \odot \pi_2 =
\begin{cases}
\text{the normalization of } \pi_1 \compose \pi_2 & \text{if } \pi_1 \compose \pi_2 \text{ is defined},\\
0 & \text{otherwise.}
\end{cases}
\]
The identity element is $\id = (\, -- \,)$. This multiplication models concatenation of linear program fragments: if $\alpha$ and $\beta$ are linear sequences of words, then
\[
\sem{\alpha\beta} = \sem{\alpha} \odot \sem{\beta}.
\]
\end{definition}

Viewed this way, $\odot$ is the typed version of the product used in earlier algebraic presentations of stack effects \cite{poial1994algebraic,poial2006typing}. It packages sequencing as a total operation with a clash element, while the checker itself continues to use the underlying partial composition $\compose$ because that formulation matches error reporting and local reasoning. The totalized view is useful mainly as a compact algebraic summary of the linear fragment; the metatheoretic results used later rely only on the operational definition of composition above.

\section{Branching and Greatest Lower Bounds of Effects}
Sequential composition is not enough for control flow. A branch requires a single summary effect that safely represents multiple alternatives.

For branching, the checker uses a partial greatest-lower-bound operation on whole effects. Intuitively, it merges branch summaries pointwise after accounting for stack cells that are preserved unchanged across both input and output. This is necessary because different branches may manipulate different depths of the stack while still agreeing on a common surrounding context.

More concretely, if one effect has $k$ more inputs than another, it must also have $k$ more outputs. The implementation then extracts a preserved prefix of length $k$ from the larger effect by requiring each corresponding input/output pair in that prefix to be comparable. After that prefix has been merged, the remaining active slices are unified pointwise. If any comparison fails, the branch summaries are incompatible.

\begin{proposition}[Whole-effect greatest lower bound]
\label{prop:effect-glb-meet}
Whenever $\pi_1 \glb \pi_2$ is defined as a whole effect:
\begin{enumerate}
\item $\pi_1 \glb \pi_2 \preceq \pi_1$,
\item $\pi_1 \glb \pi_2 \preceq \pi_2$,
\item for every effect $\rho$, if $\rho \preceq \pi_1$ and $\rho \preceq \pi_2$, then $\rho \preceq \pi_1 \glb \pi_2$.
\end{enumerate}
Hence the operator is a meet in the information order whenever it is defined.
\end{proposition}

\begin{proof}[Proof sketch]
The preserved-prefix condition is forced by stack-depth agreement: any common lower bound must preserve exactly those unmatched surrounding cells on both sides. On the active slices, the construction compares symbols pointwise and replaces each comparable pair by its type-level meet together with the induced identity constraints. That replacement is below each branch summary because every chosen symbol is at least as informative as the corresponding symbols in both inputs. It is greatest among common lower bounds because any effect below both branches must satisfy the same prefix constraints and, position by position, cannot be more informative than the comparable meet chosen by the construction.
\end{proof}

Accordingly, the checker uses greatest lower bound rather than least upper bound: it computes the strongest summary guaranteed on all paths of a conditional, not an over-approximation of behaviors that may occur on some path \cite{poial2006typing}.

This operator is used by the \texttt{either} constructor in declarative control semantics. The example
\begin{lstlisting}
: MYCHOICE IF ! ELSE C! FI ;
\end{lstlisting}
is inferred as
\begin{lstlisting}
MYCHOICE ( x a-addr flag -- )
\end{lstlisting}
because the two branches have effects
\[
(x\;a\text{-}addr -- )
\qquad\text{and}\qquad
(x\;c\text{-}addr -- ),
\]
and
\[
a\text{-}addr \glb c\text{-}addr = a\text{-}addr.
\]
The merged effect therefore preserves the more specific address requirement.

\section{Iteration and One-Step Stabilization}
The current checker does not attempt a full abstract semantics of arbitrary iteration. It does not construct a global control-flow graph or run a general fixpoint engine \cite{cousot1977abstract}. Instead, the \texttt{repeat} constructor uses a local repeated-effect operator that compares one execution of a body with two executions. This is the exact loop-related contract implemented by the checker.

\begin{definition}[One-step repeated summary]
For an effect $\pi$, define
\[
\pi^\star = \pi \glb (\pi \compose \pi),
\]
when both the composition and the greatest lower bound are defined.
\end{definition}

Operationally, this is a one-step stabilization test. The effect $\pi$ summarizes one execution of the repeated body, while $\pi \compose \pi$ summarizes two executions. If those two summaries admit a common lower bound, the checker keeps that lower bound as the information guaranteed by both one round and one further round of the body. If no such lower bound exists, the checker declines to assign a repeated-summary effect.

\begin{proposition}[Meaning of the repeated-summary operator]
Whenever $\pi^\star$ is defined:
\begin{enumerate}
\item $\pi^\star \preceq \pi$,
\item $\pi^\star \preceq \pi \compose \pi$,
\item for every effect $\rho$, if $\rho \preceq \pi$ and $\rho \preceq \pi \compose \pi$, then $\rho \preceq \pi^\star$.
\end{enumerate}
Hence $\pi^\star$ is exactly the strongest guarantee common to one and two executions of $\pi$.
\end{proposition}

\begin{proof}
By Proposition~\ref{prop:effect-glb-meet}, whole-effect $\glb$ is a meet in the information order whenever it is defined. The three clauses are therefore immediate from the definition
\[
\pi^\star = \pi \glb (\pi \compose \pi).
\]
\end{proof}

In control specifications, the \texttt{repeat} constructor uses this operator. The intended reading is not full Kleene closure and does not, by itself, claim correctness for all finite iteration counts. Rather, \texttt{repeat} accepts exactly when the one-step and two-step summaries admit a common guaranteed effect, and it returns the strongest such common effect.

The bundled program
\begin{lstlisting}
: MYLOOP BEGIN SWAP OVER WHILE INVERT REPEAT ;
\end{lstlisting}
is inferred as
\begin{lstlisting}
MYLOOP ( flag[1] flag[1] -- flag[1] flag[1] )
\end{lstlisting}
This summary is the strongest guarantee shared by one and two passes of the repeated body, which is exactly the acceptance criterion implemented for \texttt{repeat}.

\section{Declarative Specifications}
\subsection{Ordinary words, literals, and parser words}
The \texttt{specs} file associates stack effects with words and adds parsing metadata. Representative examples are:
\begin{lstlisting}
OVER ( x[2] x[1] -- x[2] x[1] x[2] )
@ ( a-addr -- x )
CONSTANT parse word define ( x -- )
VARIABLE parse word define ( -- a-addr )
literal integer ( -- n )
\end{lstlisting}

Additional clauses refine the operational role of a word.
\begin{itemize}
\item \texttt{parse word} means the word consumes the next token from the source.
\item \texttt{parse until DELIM} means the word consumes source text until a delimiter.
\item \texttt{define}, \texttt{define colon}, \texttt{define constant}, and \texttt{define variable} mark defining words.
\item \texttt{control ROLE} assigns a control-flow role such as \texttt{if}, \texttt{else}, \texttt{fi}, \texttt{begin}, or \texttt{while}.
\item \texttt{immediate}, \texttt{state interpret}, and \texttt{state compile} constrain where the word is legal.
\end{itemize}

Thus the checker is not only typing tokens; it also models the source-level behavior needed to parse Forth-like syntax in a way that aligns with standard Forth parsing disciplines \cite{forth2012standard}.

\subsection{Control structures}
Control structures are specified with paired \texttt{syntax:} and \texttt{effect:} blocks. This gives a declarative account of source-level control words, a topic that has also been studied in formal models of Forth control structures \cite{power2004formal}. For example, the sample repository contains:
\begin{lstlisting}
syntax: IF <then branch> [ELSE <else branch>] FI
  effect:
    IF
    either <then branch> <else branch>
\end{lstlisting}

and
\begin{lstlisting}
syntax: BEGIN <loop prefix> WHILE <loop body> REPEAT
  effect:
    repeat <loop prefix> WHILE
    repeat <loop body>
\end{lstlisting}

The \texttt{syntax:} clause describes a control-language pattern using concrete boundary words and metavariables for internal segments. Optional parts are written with square brackets. The \texttt{effect:} block is then compiled into an expression language with the grammar
\[
E ::= \eps \mid \langle s \rangle \mid c \mid E\,E \mid \textsf{either}(E,E) \mid \textsf{repeat}(E),
\]
where $\langle s \rangle$ denotes a named program segment and $c$ denotes the runtime effect of a control word such as \texttt{IF}, \texttt{WHILE}, or \texttt{DO}.

The denotational semantics is:
\begin{align*}
\sem{\eps} &= \id,\\
\sem{\langle s \rangle} &= \text{the inferred effect of segment } s,\\
\sem{c} &= \text{the declared runtime effect of control role } c,\\
\sem{E_1 E_2} &= \sem{E_1} \compose \sem{E_2},\\
\sem{\textsf{either}(E_1,E_2)} &= \sem{E_1} \glb \sem{E_2},\\
\sem{\textsf{repeat}(E)} &= \sem{E}^\star.
\end{align*}

This arrangement separates parsing from typing. Once a source fragment has been segmented according to a \texttt{syntax:} declaration, its meaning is calculated entirely by the effect algebra above.

\subsection{Compositional characterization}
\label{sec:constraint-view}
The declarative presentation is best read first as a system of equations over the parsed source tree. This is the exact formulation matched by the current checker. Fix a finite parsed source tree $T$ for one already-segmented source body under a finite \texttt{types}/\texttt{specs} environment $\Gamma$, and assume that every leaf token has a fixed declared summary effect in $\Gamma$. Each parsed segment $A$---a leaf token, a linear fragment, a branch arm, a loop body, or a whole definition body---is assigned an unknown summary effect $X_A$ ranging over normalized non-clashing effects together with a clash element $0$.

Write $\overline{\compose}$, $\overline{\glb}$, and $\overline{\star}$ for the totalizations of $\compose$, effect-level $\glb$, and $(\cdot)^\star$: whenever the underlying partial operation is undefined, the totalized operation returns $0$, and any occurrence of $0$ in an argument propagates to $0$. For a control template $E$, define its totalized denotation by
\begin{align*}
\widehat{\sem{\eps}}_X &= \id,\\
\widehat{\sem{\langle s \rangle}}_X &= X_s,\\
\widehat{\sem{c}}_X &= \text{the declared runtime effect of control role } c,\\
\widehat{\sem{E_1 E_2}}_X &= \overline{\compose}(\widehat{\sem{E_1}}_X,\widehat{\sem{E_2}}_X),\\
\widehat{\sem{\textsf{either}(E_1,E_2)}}_X &=
\overline{\glb}(\widehat{\sem{E_1}}_X,\widehat{\sem{E_2}}_X),\\
\widehat{\sem{\textsf{repeat}(E)}}_X &= \overline{\star}(\widehat{\sem{E}}_X).
\end{align*}

Segment equations are then generated bottom-up from $T$.
\begin{enumerate}
\item A leaf token with declared effect $\pi$ contributes $X_A = \pi$.
\item A sequential node with children $B$ and $C$ contributes
\[
X_A = \overline{\compose}(X_B,X_C).
\]
\item A control node whose \texttt{effect:} template is $E$ contributes
\[
X_A = \widehat{\sem{E}}_X,
\]
where each segment reference $\langle B \rangle$ in $E$ is interpreted as the corresponding child unknown $X_B$.
\end{enumerate}
Let $S(T,\Gamma)$ denote the resulting finite system of segment equations. Because every right-hand side refers only to proper descendants of $A$, the dependency graph of $S(T,\Gamma)$ is acyclic.

\begin{theorem}[Exactness of segment equations]
Assume that the parsed source tree $T$ is finite and that every leaf token in $T$ has a fixed declared effect under $\Gamma$. Then:
\begin{enumerate}
\item The acyclic system $S(T,\Gamma)$ has a unique solution $X^\ast$ in normalized effects plus $0$.
\item For every parsed segment $A$, the checker returns a clash exactly when $X^\ast_A = 0$; otherwise it returns the normalized effect $X^\ast_A$.
\item In particular, the checker accepts the whole parsed body if and only if the root summary satisfies $X^\ast_{\mathsf{root}} \neq 0$.
\end{enumerate}
\end{theorem}

\begin{proof}[Proof sketch]
Proceed by structural induction on the height of a node in $T$. A leaf node has a unique assigned effect by its declaration. For a sequential node or control node, the induction hypothesis gives unique summaries for all proper descendants, and the corresponding parent equation then determines a unique right-hand side because the totalized operators $\overline{\compose}$, $\overline{\glb}$, and $\overline{\star}$ are deterministic on already determined arguments. Hence $S(T,\Gamma)$ has a unique solution.

The checker evaluates parsed segments by exactly the same recursion: it first computes the summaries of child segments, then applies boundary composition for sequences, whole-effect greatest lower bound for \texttt{either}, and one-step stabilization for \texttt{repeat}. Therefore the checker's result at each node is exactly $X^\ast_A$. A clash is reported precisely when one of the partial operations used by the untotalized semantics is undefined, which is exactly the case in which the corresponding totalized equation yields $0$. The statement for the root node follows immediately.
\end{proof}

The older FORML-style inequality view can still be recovered as a relaxation by replacing each equation $X_A = F_A(\vec{X})$ with an inequality $X_A \preceq F_A(\vec{X})$ \cite{poial1991algebraic,poial1992multiple}. Every successful checker run yields a non-clashing solution of that weaker system. For the present checker, however, the exact theorem is the acyclic equation characterization above rather than a least-fixpoint result over inequalities.

\paragraph{Scope of the theorem.}
Theorem~1 is intentionally local. It characterizes one already parsed body once the summaries of all referenced leaves are fixed. The full top-level checker does more than that local semantics alone: it grows the specification dictionary as definitions are encountered, may seed provisional summaries for forward references, and supports a dedicated \texttt{RECURSE} path during definition checking. Those implementation conveniences are used to make whole-source checking practical, but they are not part of the statement of Theorem~1. The theorem covers the core body-evaluation phase that those outer mechanisms repeatedly invoke.

\section{Inference of User-Defined Words}
At top level, the checker interprets defining words immediately and uses the local body semantics of Theorem~1 to grow the specification dictionary. A colon definition starts a new sequence, parses until its terminating control word, computes the effect of the body, and installs the resulting effect as the specification of the new word. Constants and variables are handled similarly but with special defining shapes:
\[
\texttt{CONSTANT} : (x --), \qquad \texttt{VARIABLE} : (-- a\text{-}addr).
\]

The implementation therefore supports the standard Forth workflow in which the specification dictionary grows during analysis \cite{forth2012standard}. Inferred definitions are written to a log file; for the sample program the checker records:
\begin{lstlisting}
MYLOOP   ( flag[1] flag[1] -- flag[1] flag[1] )
MYCHOICE ( x a-addr flag -- )
MYTEST   ( a-addr[2] x[1] flag -- a-addr[2] x[1] a-addr[2] )
MYFI     ( x[1] x[1] flag -- x[1] x[1] )
MYAGAIN  ( a-addr[1] -- a-addr[1] )
MYSEQ    ( a-addr[2] x[1] -- x[1] a-addr[2] a-addr[2] )
MYSEQ2   ( a-addr[1] a-addr[1] -- )
\end{lstlisting}

These examples demonstrate several features of the checker at once.
\begin{itemize}
\item \texttt{MYCHOICE} uses branch merging over comparable address types.
\item \texttt{MYTEST} shows prefix-aware branch merging when the alternatives manipulate different active depths of the stack.
\item \texttt{MYLOOP} and \texttt{MYAGAIN} show loop summaries computed through repeated-effect stabilization.
\item \texttt{MYSEQ2} uses the polymorphic example word \texttt{PLUS}, showing that indexed symbols propagate through a larger composition.
\end{itemize}

\section{Relation to Prior Work}
This paper sits at the intersection of several lines of work on stack-oriented typing. Early algebraic treatments of Forth stack effects studied sequential composition, multiple stack effects, and solvability via systems of inequalities \cite{poial1991algebraic,poial1992multiple,poial1994algebraic}. Typed wildcards and subtype-sensitive matching then provided a way to track both abstract type precision and symbolic identity across stack shuffles \cite{stoddart1993type,poial2002wildcards}. Later work described branch merging and loop summaries in must-analysis terms, using greatest lower bounds and repeated-effect stabilization criteria \cite{poial2006typing}. These papers supply most of the individual ingredients used here.

\begin{table}[t]
\centering
\caption{How the present paper differs from the closest cited lines of work.}
\label{tab:prior-comparison}
\small
\begin{tabular}{@{}>{\raggedright\arraybackslash}p{0.20\linewidth}>{\raggedright\arraybackslash}p{0.26\linewidth}>{\raggedright\arraybackslash}p{0.44\linewidth}@{}}
\hline
Line of work & Main emphasis there & Difference here \\
\hline
Algebraic stack-effect specifications \cite{poial1991algebraic,poial1992multiple,poial1994algebraic} & Sequential composition, alternative effects, and inequality-style solvability & Integrates these operations into a source-level checker with declarative parser and control specifications, and replaces the global inequality claim by an exact parsed-tree characterization of what the checker computes \\
Typed wildcard and inheritance effects \cite{stoddart1993type,poial2002wildcards} & Symbolic identity tracking and subtype-sensitive polymorphism & Uses these symbolic effects as the executable summary language for whole parsed bodies, user-defined-word inference, and source-level diagnostics \\
Must-analysis typing tools \cite{poial2006typing} & Branch greatest lower bounds and loop stabilization for stack languages & Externalizes control syntax through \texttt{syntax:}/\texttt{effect:} declarations and proves exactness for the implemented parsed-body evaluator \\
Compositional low-level stack typing \cite{saabas2006compositional} & Compositional typing judgments for stack-based low-level code & Works at Forth source level, where parser words, defining words, and retargetable vocabularies are first-class inputs rather than fixed instruction forms \\
WebAssembly validation and stack-polymorphic typing \cite{webassembly2026spec,mcdermott2022wasm} & Validation of structured stack-machine bytecode with fixed control instructions and stack polymorphism around unreachable code and branches & Stays at source level rather than bytecode, treats control syntax as externally declared input, and tracks symbolic identities plus source spans for inferred word summaries and diagnostics \\
Formal control semantics \cite{power2004formal} & Semantic account of Forth control words & Targets static checking, inference, and diagnostics rather than general operational semantics \\
\hline
\end{tabular}
\end{table}

Table~\ref{tab:prior-comparison} makes the comparison explicit. The present paper is concerned mainly with how the earlier ingredients can be combined in a single checker for whole source programs. Relative to the algebraic and wildcard-based Forth papers, the novelty is not a new primitive operator but a new level of integration: declarative control specifications, parser behavior, and user-defined-word inference all live in one source-level implementation. Relative to the tooling-oriented must-analysis paper \cite{poial2006typing}, the main additional claim is an exact theorem for parsed-body checking rather than only a description of useful operators. Relative to \cite{saabas2006compositional}, the overlap is compositional reasoning for stack-oriented code, but the present work remains in the source-language idiom of Forth stack effects and parser disciplines \cite{forth2012standard}, where control syntax itself is part of the specification interface.

The WebAssembly comparison is useful for a broader programming-languages audience. Modern Wasm validation also assigns stack-machine typings compositionally and uses stack polymorphism to account for structured control flow \cite{webassembly2026spec,mcdermott2022wasm}. However, the setting is importantly different: WebAssembly validates a fixed bytecode language with built-in structured control operators, whereas the present checker analyzes source text whose parser words, defining words, and control forms are themselves supplied by external specifications. That is why this paper emphasizes declarative \texttt{syntax:}/\texttt{effect:} descriptions and source-aware diagnostics rather than only bytecode validation rules.

Relative to \cite{power2004formal}, the goal is static source checking rather than a full semantic model of control-word execution.

\section{Evaluation}
We evaluate the artifact along five axes: regression coverage, diagnostic quality, retargetability, scale, and known limitations. The reference implementation discussed throughout the paper is the Gforth checker, and the artifact additionally includes an independent Java reimplementation that consumes the same \texttt{types}/\texttt{specs}/\texttt{prog} formats. All quantitative results reported below were obtained by batch-running the bundled corpus through the Java implementation.

\paragraph{Setup and corpus.}
The bundled corpus contains 29 programs. There are 13 positive programs expected to type-check successfully and 16 negative programs expected to produce diagnostics. The positive set consists of 7 programs under the bundled Forth 2012 profile, 5 programs under custom type/specification profiles, and 1 article-sized example profile used in this paper. Table~\ref{tab:corpus-results} summarizes the current batch results.

\begin{table}[t]
\centering
\caption{Bundled regression corpus and Java batch results.}
\label{tab:corpus-results}
\begin{tabular}{lrr}
\hline
Category & Programs & Java outcome \\
\hline
Forth 2012 positives & 7 & 7 accepted \\
Custom-profile positives & 5 & 5 accepted \\
Article profile & 1 & 1 accepted \\
Negative regressions & 16 & 14 rejected \\
\hline
Total & 29 & 27 as expected \\
\hline
\end{tabular}
\end{table}

\paragraph{Regression results.}
The Java implementation accepts all 13 positive programs and rejects 14 of the 16 negative programs as expected. The two remaining negative cases are \texttt{linear\_clash.fs} and \texttt{top\_level\_clash.fs}, both of which consist of linear numeric/address clashes involving \texttt{DP 1+}. These cases expose current edge cases in top-level linear reasoning rather than failures of the declarative control machinery developed in Sections~5--6. Even so, they are valuable regression tests because they identify precisely where the current implementation is still incomplete.

\paragraph{Diagnostics.}
The negative corpus covers several distinct failure modes: malformed control nesting (\texttt{missing\_fi}, \texttt{missing\_repeat}, \texttt{unexpected\_else}), unknown words, branch-summary incompatibility, loop-body incompatibility, nested-definition rejection, unterminated definitions, and multiple-error recovery. Diagnostics include file positions and the offending source span. For example, on \texttt{if\_clash.fs} the checker reports that the two \texttt{IF} alternatives have effects \texttt{( -- a-addr )} and \texttt{( -- char )} and highlights the exact \texttt{IF ... ELSE ... FI} region. This is the level of feedback needed to debug a declarative stack-effect clash rather than merely observe that checking failed.

\paragraph{Retargetability.}
The custom-profile tests exercise features that are not hard-coded into the checker core: custom literal handling, custom type hierarchies, aliasing profiles, and declaratively specified control structures. In particular, the bundled \texttt{WHEN ... OTHERWISE ... ENDWHEN} example and the switch-like example in \texttt{positive/custom\_types/} are described through \texttt{syntax:}/\texttt{effect:} declarations rather than evaluator-specific control code. Together with the literal-alias, literal-count, and literal-scanner examples, these tests show that the same effect calculus can be reused across profiles with substantially different vocabularies.

\paragraph{Scale and running time.}
Table~\ref{tab:artifact-scale} summarizes the main scale indicators in the artifact. The repository's bundled Forth 2012 environment consists of a 50-line type file and a 278-line specification file, for 328 lines of declarative typing and control information in total. Under that environment, the checker analyzes the single-file implementation \texttt{gforth-evaluator.fs}, which has 4{,}255 lines. The artifact also contains an independent Java implementation totaling 7{,}555 lines across its \texttt{.java} sources, giving a second executable realization of the same input language. On one development machine, the Java runtime analyzes a representative bundled Forth 2012 demo in about 0.13 seconds, the article profile in about 0.07 seconds, and the full 29-program regression corpus in about 3.16 seconds.

\begin{table}[t]
\centering
\caption{Representative artifact scale and Java running times.}
\label{tab:artifact-scale}
\begin{tabular}{lr}
\hline
Item & Size / time \\
\hline
Forth 2012 type declarations & 50 lines \\
Forth 2012 specifications & 278 lines \\
Bundled analyzed application & 4{,}255 lines \\
Independent Java implementation & 7{,}555 LOC \\
Single scanner demo & 0.13 s \\
Article profile & 0.07 s \\
Full 29-program corpus & 3.16 s \\
\hline
\end{tabular}
\end{table}

\paragraph{Limitations and scope.}
The checker is designed as a static evaluator rather than a full runtime. It relies on the declared specification set, rejects nested defining words inside definitions, enforces interpretation-state restrictions, and treats the external type hierarchy as the semantic authority. The regression corpus is therefore best read as evidence of coverage for the intended source-language fragment, not as a claim of broad industrial benchmarking. Within that scope, however, the evaluation shows that the same checker core supports both the paper's theory examples and larger bundled artifacts, while also surfacing two concrete remaining precision gaps.

\section{Conclusion}
The paper makes two main points. First, typed stack-effect checking for a Forth(-like) language can be organized around a small collection of operations on symbolic effects: sequential composition, branch merging by greatest lower bound, and a local loop-summary operator for repeated control flow. Second, for fixed parsed bodies under finite specifications and fixed leaf summaries, this operational checker has an exact compositional characterization over the parsed source tree: acceptance coincides with the absence of clashes in the induced segment equations, and a successful run yields their unique non-clashing summaries.

Because the type hierarchy and the word and control specifications are external to the core algorithm, the same implementation can be retargeted to different vocabularies and dialects with relatively little change to the checker itself. The evaluation shows this on three scales at once: a 29-program regression corpus, several custom control and literal profiles, and the larger bundled Forth 2012 artifact discussed above \cite{forth2012standard}. The broader methodological point is that many source-level features of a concatenative language become amenable to static reasoning once they are phrased as operations in a finite effect domain.

\bibliographystyle{ACM-Reference-Format}
\bibliography{typechecker-paper-popl-public}

\end{document}
